\documentclass[11pt,a4paper]{article}

% Working notes for the Neural Particle Method project.
% Synthetic-data setup, machinery, and the benchmark suite of 9 September 2026.

\usepackage{amsmath, amssymb, amsthm, mathtools}
\usepackage{booktabs}
\usepackage{graphicx}
\usepackage{array}
\usepackage[numbers,sort&compress]{natbib}
\usepackage[colorlinks=true,linkcolor=blue,citecolor=blue,urlcolor=blue]{hyperref}
\usepackage[capitalise,noabbrev]{cleveref}
\usepackage{microtype}
\usepackage{parskip}
\setlength{\parindent}{0pt}
\setlength{\parskip}{6pt}
\usepackage[margin=2.2cm]{geometry}
\usepackage{xcolor}
\setlength{\emergencystretch}{3em}

\newcommand{\E}{\mathbb{E}}
\newcommand{\R}{\mathbb{R}}
\newcommand{\sigDup}{\sigma_{\mathrm{Dup}}}
\newcommand{\Lev}{L}
\newcommand{\Lmax}{L_{\max}}
\newcommand{\bp}{\,\mathrm{bp}}
\newcommand{\code}[1]{\texttt{#1}}

\title{Neural Particle Method: Experimental Record\\ \large Synthetic-data setup, machinery, and the benchmark suite}
\author{Osian Shelley}
\date{8 September 2026}

\begin{document}
\maketitle

\begin{abstract}
This note records, in one place, the synthetic markets and dynamics used throughout the project, the simulation, estimation and scoring machinery, and the specification of the clean benchmark suite agreed on 8 September 2026 (\cref{sec:suite}), run on 9 September 2026; its tables are generated from the MLflow store. The suite follows the discovery, in the diagnostics of 8 September, that all one-pass schemes were fitting the conditional expectation of the raw Euler variance rather than of its positive part, which the dynamics actually use (\cref{sec:floor}); the diagnostic runs themselves live in the MLflow experiments \code{mechanism}, \code{bump\_floor} and \code{bump\_floor\_nn}.
\end{abstract}

\tableofcontents

\section{Conventions}\label{sec:conv}

\emph{Explicit} denotes the one-pass, per-slice calibration with a per-slice neural network. Nadaraya--Watson (NW) is the classical particle method: the same one-pass structure with kernel regression. Other per-slice estimators are named by their estimator (RKHS, spline, ridge). \emph{Implicit} denotes the damped outer iteration with a single network in $(t,x)$. All errors are implied-volatility root-mean-square errors in volatility basis points against the target surface, unless stated otherwise. ``Pooled'' means over all quoted strikes and maturities together; ``wings'' means the quoted strikes with $|k|>0.25$; per-maturity columns restrict to one maturity.

\section{Synthetic data}\label{sec:data}

\subsection{Target surfaces: SSVI power law}\label{sec:ssvi}

Each scenario's market is an SSVI implied-variance surface in the power-law parameterisation. With log-moneyness $k=\log(K/S_0)$, $S_0=1$, and parameters $(\sigma_0,\eta,\gamma,\rho_s)$,
\begin{equation}\label{eq:ssvi}
  \theta_T = \sigma_0^2 T,\qquad
  \varphi(\theta_T) = \eta\,\theta_T^{-\gamma},\qquad
  w(k,T) = \tfrac12\theta_T\Bigl(1+\rho_s\varphi k+\sqrt{(\varphi k+\rho_s)^2+1-\rho_s^2}\Bigr),
\end{equation}
where $w$ is total implied variance and the implied volatility is $\sqrt{w(k,T)/T}$. Parameters are drawn independently per scenario from
\[
  \sigma_0\sim U(0.12,0.35),\quad \eta\sim U(0.5,2.0),\quad \gamma\sim U(0.3,0.5),\quad \rho_s\sim U(-0.85,-0.10),
\]
resampled until the sufficient no-arbitrage conditions $0<\gamma\le\tfrac12$, $|\rho_s|<1$, $\eta(1+|\rho_s|)\le 2$ hold. Scenario $s_i$ uses the seed $1000+i$ (\code{bench/scenarios.py}, \code{make\_registry}).

The Dupire local volatility is analytic from \cref{eq:ssvi}: with $w_k,w_{kk},w_T$ the derivatives of $w$,
\begin{equation}\label{eq:dupire}
  \sigDup^2(T,k) = \frac{w_T}{1-\frac{k}{w}w_k+\frac14\bigl(-\frac14-\frac1w+\frac{k^2}{w^2}\bigr)w_k^2+\frac12 w_{kk}},
\end{equation}
with the denominator floored at $0.05$, the local variance clipped to $[10^{-8},9]$, and the time argument clipped to $[t_{\min},T]$ with $t_{\min}=0.004$ (\code{market/local\_vol.py}). These clipping constants are part of the calibrated behaviour and are never changed.

\subsection{Variance dynamics}\label{sec:dyn}

The calibrated model is Heston-type LSV in log-spot $X=\log S$,
\begin{equation}\label{eq:lsv}
  dX_t = -\tfrac12\Lev^2(t,X_t)V_t\,dt+\Lev(t,X_t)\sqrt{V_t}\,dW_t,\qquad
  dV_t = \kappa(\theta-V_t)\,dt+\xi\sqrt{V_t}\,dB_t,\qquad d\langle W,B\rangle_t=\rho\,dt .
\end{equation}
For scenario $s_i$ the dynamics parameters are drawn from the same seed as the surface: $\kappa\in\{1,2,3\}$, $\xi\in\{0.3,0.6,1.0\}$, $\rho\in\{-0.7,-0.3\}$, each uniform over its set, and $\theta=v_0=\sigma_0^2$ so that the variance level matches the surface's short-end level. The Feller ratio $2\kappa\theta/\xi^2$ therefore ranges from $0.05$ to $4.3$ across the grid. \Cref{tab:scen} lists all twenty scenarios. Horizon $T=2$ years; spot $S_0=1$; no rates or dividends.

\begin{table}[htbp]
\centering\small
\begin{tabular}{l rrrr rrr r}
\toprule
 & \multicolumn{4}{c}{SSVI surface} & \multicolumn{3}{c}{dynamics} & \\
\cmidrule(lr){2-5}\cmidrule(lr){6-8}
sid & $\sigma_0$ & $\eta$ & $\gamma$ & $\rho_s$ & $\kappa$ & $\xi$ & $\rho$ & Feller\\
\midrule
s01 & 0.261 & 0.52 & 0.34 & $-0.21$ & 1 & 0.3 & $-0.7$ & 1.51 \\
s02 & 0.208 & 1.04 & 0.45 & $-0.56$ & 3 & 0.6 & $-0.7$ & 0.72 \\
s03 & 0.163 & 0.89 & 0.40 & $-0.35$ & 3 & 1.0 & $-0.7$ & 0.16 \\
s04 & 0.120 & 0.70 & 0.33 & $-0.56$ & 2 & 0.6 & $-0.3$ & 0.16 \\
s05 & 0.285 & 1.01 & 0.49 & $-0.72$ & 2 & 1.0 & $-0.3$ & 0.32 \\
s06 & 0.304 & 1.00 & 0.38 & $-0.58$ & 3 & 1.0 & $-0.3$ & 0.55 \\
s07 & 0.202 & 0.56 & 0.46 & $-0.49$ & 3 & 1.0 & $-0.3$ & 0.25 \\
s08 & 0.158 & 1.72 & 0.34 & $-0.11$ & 1 & 1.0 & $-0.7$ & 0.05 \\
s09 & 0.133 & 0.99 & 0.31 & $-0.64$ & 3 & 1.0 & $-0.3$ & 0.11 \\
s10 & 0.142 & 1.30 & 0.48 & $-0.11$ & 3 & 1.0 & $-0.7$ & 0.12 \\
s11 & 0.144 & 1.10 & 0.38 & $-0.11$ & 2 & 1.0 & $-0.7$ & 0.08 \\
s12 & 0.191 & 1.01 & 0.33 & $-0.22$ & 3 & 1.0 & $-0.7$ & 0.22 \\
s13 & 0.303 & 0.62 & 0.43 & $-0.51$ & 2 & 1.0 & $-0.3$ & 0.37 \\
s14 & 0.220 & 1.38 & 0.48 & $-0.11$ & 1 & 0.6 & $-0.3$ & 0.27 \\
s15 & 0.127 & 0.72 & 0.37 & $-0.33$ & 3 & 0.3 & $-0.3$ & 1.08 \\
s16 & 0.172 & 1.16 & 0.37 & $-0.42$ & 1 & 0.6 & $-0.7$ & 0.17 \\
s17 & 0.136 & 1.13 & 0.42 & $-0.50$ & 1 & 0.3 & $-0.3$ & 0.41 \\
s18 & 0.257 & 1.18 & 0.47 & $-0.24$ & 3 & 0.6 & $-0.3$ & 1.10 \\
s19 & 0.170 & 1.09 & 0.39 & $-0.23$ & 2 & 0.3 & $-0.7$ & 1.28 \\
s20 & 0.253 & 0.93 & 0.32 & $-0.32$ & 3 & 0.3 & $-0.7$ & 4.27 \\
\bottomrule
\end{tabular}
\caption{The twenty synthetic scenarios: SSVI surface parameters, variance dynamics, and the Feller ratio $2\kappa\theta/\xi^2$.}
\label{tab:scen}
\end{table}

Two further families exist in the registry and were not used today: six variants of s01 with $(\xi,\rho)$ swept over the grid (\code{fig3\_registry}), and three Heston-market scenarios whose target surface is generated by a Heston model with parameters different from the calibrated dynamics (\code{li\_simple}, \code{li\_complex}, \code{bayer}; horizon 1 year; Dupire surface from prices on a $50\times121$ grid).

\subsection{Quotes, bump, and scoring}\label{sec:quotes}

Quoted strikes are $13$ log-strikes equally spaced in log between $0.6$ and $1.6$; maturities are $0.25$, $0.5$, $1$ and $2$ years. The wings are the strikes with $|k|>0.25$ (seven of thirteen). The warm-start experiments bump the surface to $\sigma_0'=\sigma_0+0.01$, $\eta'=0.95\eta$, $\rho_s'=\min(\rho_s+0.03,-0.05)$, a level rise of one volatility point with curvature and skew nudged inside the no-arbitrage region.

Scoring reprices vanilla calls under the calibrated leverage field on fresh paths: $500\,000$ particles, $200$ Euler steps over $[0,2]$, seed $10\,000$ plus a strategy-specific offset, call price as the sample mean of the payoff, Black--Scholes inversion at the grid-snapped maturity, and the exact SSVI implied volatility at the same snapped time as the target (\code{pricing/reprice.py}, \code{pricing/metrics.py}). The bump experiments use $300\,000$ repricing paths. Errors are reported as RMSE in volatility bp, pooled, wings-only, and per maturity.

\section{Machinery}\label{sec:machinery}

\subsection{Simulation}\label{sec:sim}

Calibration and repricing share one stepper. With $\Delta=T/K$, $v^+=\max(v,0)$, $z_1=\rho z_b+\sqrt{1-\rho^2}z_p$ and independent standard normals $z_b,z_p$,
\begin{equation}\label{eq:euler}
  x_{k+1} = x_k-\tfrac12\Lev_k(x_k)^2v_k^+\Delta+\Lev_k(x_k)\sqrt{v_k^+}\sqrt{\Delta}\,z_1,\qquad
  v_{k+1} = v_k+\kappa(\theta-v_k^+)\Delta+\xi\sqrt{v_k^+}\sqrt{\Delta}\,z_b .
\end{equation}
This is full truncation: the raw $v$ may be negative and is carried, but only $v^+$ enters any coefficient. The paper calibrated with $K=50$ steps; the benchmark suite of \cref{sec:suite} uses $200$ for calibration and repricing alike. The leverage is stored per slice on a grid in log-spot and evaluated by linear interpolation, constant beyond the grid ends; the slice in force at time $t$ is the last grid time $\le t$. The explicit and NW passes use a per-slice grid at the $0.1\%$ to $99.9\%$ quantiles of the current cloud ($101$ points); the implicit loop uses a fixed grid of $81$ points from $\log 0.4$ to $\log 2.2$. The leverage is capped at $\Lmax=4$ and the fitted conditional expectation is floored at $10^{-4}$.

\subsection{One-pass schemes and their estimators}\label{sec:onepass}

At slice $k$, fit $f_k(x)\approx\E[V_{t_k}\mid X_{t_k}=x]$ on a subsample of $30\,000$ particles from the current cloud (the earlier default of the code; the benchmark suite of \cref{sec:suite} fits on the full cloud), set $\Lev_k=\operatorname{clip}(\sigDup(t_k,\cdot)/\sqrt{\max(f_k,10^{-4})},0,\Lmax)$, and advance the cloud. Slice $0$ uses $f_0\equiv v_0$. Estimators:
\begin{itemize}
\item \textbf{NW}: Gaussian Nadaraya--Watson with Silverman bandwidth $h=1.06\,\hat\sigma_x\,n^{-1/5}$ on the subsample.
\item \textbf{Explicit (per-slice network)}: a $1\to64\to64\to1$ SiLU network with softplus output scaled by $0.04$, input scaled by $0.3$, trained by Adam at learning rate $10^{-2}$ on the unweighted squared loss; $400$ steps on the first fitted slice and $120$ steps on each later slice, warm-started from the previous slice's weights.
\item \textbf{RKHS}: kernel ridge with a Gaussian kernel of variance $0.1$ in log-spot, $100$ centres at the $j\cdot100/101$ percentiles of the cloud, regularisation $\lambda=10^{-9}$ (Bayer et al.'s published choices).
\item \textbf{Ridge head on a frozen body}: the implicit network's body (below) as a $65$-dimensional feature map, with a per-slice ridge readout, $\lambda=10^{-3}$ per sample, shrunk toward the previous slice's coefficients.
\end{itemize}

\subsection{Implicit loop}\label{sec:implicit}

A single network $f_\theta(t,x)$, $2\to64\to64\to1$ SiLU with softplus output scaled by $0.04$, inputs $(t/T,\,x/0.3)$. Starting from the explicit output as $\Lev^0$, each outer iteration simulates a full cloud under $\Lev^n$ with fresh random numbers, draws a pooled subsample of $60\,000$ points across all slices, trains $\theta$ for $300$ Adam steps at learning rate $10^{-2}$ with the network and optimiser state carried over from the previous iteration, reads $\Phi(\Lev^n)_k=\operatorname{clip}(\sigDup(t_k)/\sqrt{\max(f_\theta(t_k,\cdot),10^{-4})},0,\Lmax)$ on the fixed grid, and sets $\Lev^{n+1}=(1-\alpha)\Lev^n+\alpha\Phi(\Lev^n)$ with $\alpha=\tfrac12$. The paper uses six iterations. The pooled data are truncated, $v^+$, before training.

\subsection{The variance-target discrepancy}\label{sec:floor}

Before 8 September 2026 the one-pass schemes handed the estimator the raw Euler $v$, while the implicit loop handed it $v^+$. Since only $v^+$ drives the dynamics \cref{eq:euler}, the quantity the leverage must divide by is $\E[V^+\mid X]$. Under Feller violation a sizeable fraction of particles has $v<0$ after a coarse step (on s11, $v_0=0.02$, $\xi=1$, $\Delta=0.04$: the one-step standard deviation of $v$ is $0.028$), so $\E[v\mid X]<\E[v^+\mid X]$ by the mean negative excursion, the denominator was too small, and the leverage too large. The effect is independent of the particle count and shrinks with the step size. Fitting on $v^+$ (\code{ExplicitConfig.fit\_v\_floor}) has been the package default since 8 September 2026; every result in \cref{sec:suite} and the study log (\code{paper/studies\_log.md}) uses it, and the archived replays of the paper's runs pin it off explicitly so that they still reproduce bit for bit.

\section{Benchmark suite}\label{sec:suite}

This section specifies the clean benchmark agreed on 8 September 2026 and run on 9 September 2026 with the variance floor of \cref{sec:floor} as the package default. Every cell of the tables below is one MLflow run, keyed by method, scenario, seed and, for the online stage, lag type and particle budget, in experiments \code{suite\_cold}, \code{suite\_offline}, \code{suite\_lagged}, \code{suite\_budget\_tuned} and \code{pde\_reference}. Every trained object is stored as a run artifact (network weights as a \code{torch} state dictionary, per-slice estimators and leverage fields as pickles and JSON) and can be reloaded by run id. The tables are generated from the store by \code{nparticle suite tables}.

\subsection{Common setup}\label{sec:suite-setup}

\begin{itemize}
\item \textbf{Scenarios.} The $20$ SSVI scenarios of \cref{tab:scen}. The three Heston-market scenarios used in the literature comparisons (\code{li\_simple}, \code{li\_complex}, \code{bayer}; a Heston market surface, a different Heston variance process as the calibrated model, one-year maturity, log-strikes on $[-0.9,0.9]$) are run in the same experiments and are in the store, but they score a single one-year maturity and are not reported in this draft.
\item \textbf{Time grid.} $200$ Euler steps for calibration and $200$ for repricing, $\Delta=0.01$ on $[0,2]$, so that every method sits at its step-count plateau (the step-count diagnostics of 8 September 2026, MLflow experiment \code{mechanism}) and the comparison is about the estimator alone.
\item \textbf{Particles.} $100\,000$ particles for every cold row and for the online stage at its full budget; at those sizes every per-slice fit and every online head sees the whole cloud, and the earlier code default of a $30\,000$-particle fit subsample is not used. The online budget is also swept below $100\,000$, and there too every fit uses the whole online cloud. Offline bodies at $200\,000$ and $500\,000$ particles: the fit subsample stays at $100\,000$, so each per-slice fit of a body reads a fresh $100\,000$-particle subsample of the offline cloud while the cloud itself is simulated at the full size.
\item \textbf{Seeds.} Two seeds ($0$, $1$) for every cold row and every online stage; one seed for each offline body. The reprice seed is $10\,000$ plus the calibration seed.
\item \textbf{Floor.} All one-pass schemes fit $\E[V^+\mid X]$; \code{fit\_v\_floor} becomes the default in this work and the tiny golden replays are regenerated in the same commit.
\item \textbf{Implicit configuration.} Warm start from the explicit network field, $30$ damped iterations at $\alpha=\tfrac12$ with fresh noise per iteration, persistent network and optimiser (the best practical setting found in the convergence diagnostics of 8 September 2026, experiment \code{mechanism}).
\item \textbf{Scoring.} As in \cref{sec:quotes}: $500\,000$ repricing paths, Black--Scholes inversion at the snapped maturity, exact SSVI (or semi-analytic Heston) implied volatility as the target. Metrics per run: pooled and wings MAE and RMSE in volatility bp, both per maturity, the count of failed inversions, and the leverage RMSE against the PDE reference of the scenario. Errors are reported in implied-vol basis points; each table also shows the liquid-quote error (mean absolute error over $|k|\le0.25$ and $T\ge0.5$) and the absolute option-price error in basis points of spot, recovered from the implied-vol errors through the Black--Scholes formula.
\item \textbf{Latency.} Wall-clock seconds of the calibration only (repricing excluded). For the lagged suite, only the online stage is timed.
\item \textbf{PDE reference.} The Fokker--Planck reference leverage (\code{reference/pde.py}) is solved per scenario at $200$ steps and enters twice: repriced under the same Monte Carlo as every other row, which is the attainable floor of the scoring under this repricer (Gy\"ongy), and as the reference for the leverage-error column. Its repricing is repeated under a second reprice seed to measure the Monte Carlo noise floor of the score itself.
\end{itemize}

\subsection{Cold suite}\label{sec:suite-cold}

One calibration from scratch per method, scenario and seed at $100\,000$ particles. Methods, with their estimator of $\E[V^+\mid X]$ per slice:
\begin{enumerate}
\item \textbf{NW}: Gaussian Nadaraya--Watson, Silverman bandwidth (\cref{sec:onepass}).
\item \textbf{Explicit NN, short training}: per-slice network, warm-started across slices, $400$ Adam steps on the first slice and $120$ on each later slice, full batch (\cref{sec:onepass}).
\item \textbf{Explicit NN, tuned}: the same per-slice scheme with a $1\to64\to64\to64\to1$ SiLU network, minibatch Adam (batches of $8192$, learning rate $10^{-3}$), $2000$ steps on the first slice and $500$ warm-started steps on each later slice; the architecture and schedule that the study log (\code{paper/studies\_log.md}) found accurate in the wings and robust under Feller violation.
\item \textbf{Explicit NN, searched}: the per-slice scheme with the recipe found by the Optuna search recorded in the study log (\code{paper/studies\_log.md}): a $1\to32\to32\to32\to32\to1$ SiLU network, minibatch Adam (batches of $8192$, learning rate $3.7\times10^{-3}$), $500$ steps on the first slice and $400$ warm-started steps on each later slice, fitted on $250\,000$ particles of the cloud, with mean matching (the fitted slice is rescaled to reproduce the cloud's mean variance), a monotone penalty on slopes of the wrong sign, and a linear continuation of the slice beyond its quantile grid.
\item \textbf{Implicit NN}: damped global-network loop (\cref{sec:implicit}, configuration above).
\item \textbf{RKHS}: Gaussian kernel ridge with $100$ quantile centres, Bayer et al.'s published knobs (\cref{sec:onepass}).
\item \textbf{Spline}: penalised cubic B-spline, $25$ quantile knots, $\lambda=1$.
\item \textbf{Literature block.} GHL kernel: Nadaraya--Watson in spot with the quartic kernel and the Guyon--Henry-Labord\`ere rule-of-thumb bandwidth $h(t)=1.5\,S_0\,\sigma_{\mathrm{LV}}(S_0,t)\sqrt{\max(t,0.25)}\,N^{-1/5}$. Bins: piecewise-constant regression on $20$ equal-frequency buckets (van der Stoep et al.). Muguruza: the kernel-free conditional Monte Carlo of Corollary~4.1, each particle contributing its one-step Gaussian conditional law of log-spot given the variance path. PU-RBF: Hakala's partition-of-unity radial basis functions, $40$ centres, $5$-nearest-neighbour widths, regulariser $0.2$.
\item \textbf{PDE}: the reference leverage repriced by Monte Carlo (attainable floor).
\end{enumerate}

\begin{table}[h]\centering\footnotesize
\caption{Cold suite, $20$ SSVI scenarios, mean and median over scenarios and two seeds. Liquid: $|k|\le0.25$, $T\ge0.5$. The PDE row is the attainable floor of the scoring, its seed-to-seed spread is $6.6\bp$ (the Monte Carlo noise floor of the score), and its latency is the Fokker--Planck solve. Bold: the best method row in each column; the PDE floor is the reference and is not a contender.}\label{tab:suite-cold}
\begin{tabular}{l rrrrr r}
\toprule
method & \multicolumn{4}{c}{implied vol, bp} & price, bp & s\\
\cmidrule(lr){2-5}
& mean & median & liquid & wings & pooled & median\\
\midrule
NW & 46 & 33 & 16 & 70 & 2.7 & 8.0\\
Explicit NN, short training & 57 & 36 & 16 & 86 & 2.6 & 1042.7\\
Explicit NN, tuned & 49 & 35 & 14 & 75 & 2.4 & 536.3\\
Implicit NN & 58 & 37 & 19 & 88 & 3.3 & 1298.6\\
RKHS & 81 & 70 & 45 & 113 & 8.5 & 12.4\\
Spline & 47 & 34 & 14 & 73 & 2.5 & 2.7\\
GHL kernel & 45 & 32 & 15 & 70 & 2.5 & 8.5\\
Bins & 66 & 61 & 16 & 105 & 2.7 & 1.5\\
Muguruza & 154 & 139 & 184 & 147 & 29.1 & 10.6\\
PU-RBF & 57 & 40 & 16 & 90 & 2.8 & 2.5\\
PDE (attainable floor) & 42 & 27 & 13 & 62 & 2.2 & 207.3\\
\bottomrule
\end{tabular}

\end{table}

\emph{Results} (\cref{tab:suite-cold}). With the variance floor in place the classical kernels are the best cold calibrators: the GHL kernel and NW sit at $45$--$46\bp$ pooled, $3$--$4\bp$ above the attainable floor of $42$, and the spline at $47$. The tuned per-slice network closes most of the gap the short-trained network leaves ($57\to49\bp$ pooled, $86\to75$ in the wings) and is the best row on liquid quotes ($14\bp$, against $15$--$16$ for the kernels and $13$ for the floor) and in price space ($2.4\bp$ of spot against $2.5$--$2.7$), but it does not beat the kernels on the pooled score: it is ahead of NW on $12$ of the $20$ scenarios, by a couple of basis points each, and behind by $14$--$39\bp$ on s03, s08 and s11, the scenarios where the variance spends most of its time at zero. Its latency is $536$ s against $8$ for NW, so as a cold calibrator it is a match, not an improvement. The searched network (study log) is $48\bp$ pooled, between NW and the hand-tuned network, better in the wings ($72$) and worse on liquid quotes ($17$), at $895$ s: the search improved the estimator of $\E[V\mid X]$ by a few percent, and that does not change the cold picture. The implicit loop is $58\bp$ and the slowest row. Bins, PU-RBF, RKHS and Muguruza's conditional Monte Carlo are behind the kernels by $10$ to $110\bp$, RKHS and Muguruza mostly in the wings. Every row is dominated by the three-month wings; on liquid quotes all the competitive rows are within $3\bp$ of the floor, and in price space within $0.5\bp$ of spot.


\subsection{What a lag is}\label{sec:suite-lag}

An overnight calibration is done on surface $S_0$; the intraday market is $S_1$, and the question is how much of the overnight work survives. Because the synthetic surfaces are parametric, a lag is a move in parameter space and never a bump of quoted points, so arbitrage cannot be introduced by accident. Two lags are run for every cell.

\emph{Surface lag.} For SSVI, the existing bump: $\sigma_0'=\sigma_0+0.01$, $\eta'=0.95\eta$, $\rho_s'=\min(\rho_s+0.03,-0.05)$, shrunk by quarters until the Gatheral--Jacquier sufficient conditions ($\eta(1+|\rho_s|)\le2$, $0<\gamma\le\tfrac12$) hold; calendar no-arbitrage is automatic since $\theta_T=\sigma_0^2T$ is increasing. For a Heston market, the matching move in vol terms: $\sqrt{v_0}$ and $\sqrt\theta$ up one volatility point, $\rho\to\min(\rho+0.03,-0.05)$, $\xi\to0.95\xi$; the surface is model-generated and arbitrage-free for any valid parameters.

\emph{Surface plus spot lag.} The surface lag together with a sticky-strike spot move of $+2\%$, $S_1=S_0e^{\delta}$, $\delta=\log1.02$. Under sticky-strike the implied volatility at each absolute strike is unchanged, so in log-moneyness relative to the new spot the surface is translated: $\sigma_{\mathrm{IV}}'(k,T)=\sigma_{\mathrm{IV}}(k+\delta,T)$. A translation in $k$ preserves Durrleman's butterfly condition (it involves only $w$ and its $k$-derivatives at a point) and calendar monotonicity, so the shifted surface is arbitrage-free. Call prices at fixed strike change with the spot, so the Dupire local volatility is that of the translated surface, $\tilde w(k)=w(k+\delta)$ with $k$ measured from the new forward; the particles also start at the new spot. The stored denominators $f(t,x)=\E[V_t\mid X_t=x]$ were fitted for a cloud started at $S_0$, and a conditional expectation given spot depends on the starting point: with $\rho<0$ a particle sitting $2\%$ above the old start had its variance drift down, so the stale $f$ at the new spot is below $v_0$, the stale leverage is too high near the money, most at short maturities. Sticky-delta, by contrast, leaves the surface in moneyness unchanged and the whole calibration is scale invariant, so a sticky-delta spot move is invisible to every method; the surface-only lag already covers it. A $2\%$ daily move is one to one-and-a-half standard deviations at $20$--$30\%$ volatility.

\subsection{Offline bodies}\label{sec:suite-bodies}

\emph{Offline stage} (experiment \code{suite\_offline}, one run per scenario, body type and particle count): a full calibration on $S_0$ at $200\,000$ or $500\,000$ particles, $200$ steps, seed $0$. Three bodies are trained. The \textbf{explicit NN, short training} is the per-slice network pass on the paper's schedule ($400$ Adam steps on the first slice, $120$ on each later slice, full batch); the \textbf{explicit NN, tuned} is the same pass with the network and schedule of item~3 of \cref{sec:suite-cold} ($1\to64\to64\to64\to1$ SiLU, minibatch Adam, $2000$ steps on the first slice and $500$ on each later slice); both save the per-slice networks and the per-slice denominators $f_k$. The \textbf{implicit NN} is the damped global loop of \cref{sec:implicit} and saves the global network. Each body is scored here as a cold calibrator on its own overnight surface $S_0$, against the same repricer and the same quotes as \cref{sec:suite-cold}, so the table says what a body is worth before any lag. These runs are cached and reused by every online cell below.

\begin{table}[h]\centering\footnotesize
\caption{Offline bodies as cold calibrators on $S_0$, $20$ SSVI scenarios, seed $0$; training time is the full offline calibration. Bold: the best method row in each column; the PDE floor is the reference and is not a contender.}\label{tab:suite-bodies}
\begin{tabular}{l l rrr r r}
\toprule
body & particles & pooled & liquid & wings & price, bp & s\\
\midrule
Explicit NN, short training & 200k & 54 & -- & 83 & 2.6 & 1047.2\\
Explicit NN, short training & 500k & 55 & -- & 84 & 2.7 & 1049.0\\
Explicit NN, tuned & 200k & -- & -- & -- & -- & --\\
Explicit NN, tuned & 500k & -- & -- & -- & -- & --\\
Implicit NN & 200k & 59 & -- & 93 & 2.9 & 1365.5\\
Implicit NN & 500k & 62 & -- & 95 & 3.3 & 1469.2\\
PDE (attainable floor) & -- & 42 & -- & 62 & 2.2 & 207.3\\
\bottomrule
\end{tabular}

\end{table}

\emph{Results} (\cref{tab:suite-bodies}). The tuned network is the only body worth keeping: $47\bp$ pooled at either size, against $54$--$55$ for the short-trained network and $59$--$62$ for the implicit loop, and it trains in half their time. More offline particles do not help any body at these training lengths ($200$k and $500$k are within $1\bp$), and the implicit loop gets slightly worse with more particles; the fit, not the cloud, is the limit. The tuned body's liquid-quote error ($14\bp$) is within $1\bp$ of the floor. The searched body at $500\,000$ particles is $48\bp$ pooled, the best of the bodies in the wings ($70$) and the worst on liquid quotes ($18$); mean matching moves the near-money level while it fixes the wings. It trains in $856$ s. (Only the $500\,000$-particle searched body was trained; it is the one the online head uses.)


\subsection{Frozen body and online head}\label{sec:suite-heads}

\emph{Online stage} (one run per scenario, method, lag, online budget and online seed): given the intraday surface $S_1$ and a fresh cloud, build the intraday leverage and reprice. Only this stage is timed. The body is the tuned $500\,000$-particle body of \cref{sec:suite-bodies}, frozen; the methods differ in what is refitted online:
\begin{enumerate}
\item \textbf{Tuned body, stale $f$, fresh Dupire}: $\Lev_k=\operatorname{clip}(\sigDup^{S_1}(t_k,\cdot)/\sqrt{\max(f_k,10^{-4})},0,\Lmax)$ with the offline per-slice $f_k$. Costs nothing online beyond the repricing.
\item \textbf{Tuned body with spline head}: simulate the online particles under the stale-$f$ leverage of item~1 with fresh Dupire, slice by slice; at each slice fit the penalised cubic B-spline of \cref{sec:onepass} ($25$ quantile knots, $\lambda=1$) to the residual $v^+-f_k(x)$ on the online cloud, and use $\max(f_k+\text{residual},\,\tfrac14 f_k)$ as the denominator for the next step; the floor keeps the head from driving the leverage to its cap on Feller-violating scenarios.
\item \textbf{Tuned body with RKHS head}: the same residual head with a Gaussian kernel ridge (kernel width $0.5$ times the cloud's standard deviation, ridge $10^{-3}$) in place of the spline, with the same floor.
\item \textbf{Tuned body with ridge head}: the per-slice network body frozen as a feature map ($64$ features plus a constant) with its last layer refitted by ridge on the online cloud, shrunk toward the offline readout. The body's output is $\operatorname{softplus}(A w)\cdot 0.04$, and the refit minimises $\sum_i(\operatorname{softplus}(A_i w)\cdot 0.04-v_i^+)^2+\lambda_{\mathrm{eff}}\|w-w_0\|^2$ on the raw variance scale by Gauss--Newton to convergence (at most $20$ steps) from the offline last layer $w_0$, with $\lambda_{\mathrm{eff}}=\lambda$ times the mean diagonal of the Gauss--Newton curvature at $w_0$ and $\lambda=10^{-3}$, so the penalty is a fixed fraction of the data curvature whatever the output scale and cloud size. The denominator is $\operatorname{softplus}(A w)\cdot0.04$ at the converged $w$. When the online data agree with the offline fit the first step is zero and $w=w_0$.
\end{enumerate}
The online reference row is NW re-solved from scratch on $S_1$ at the same budget and $200$ steps (the best intraday strategy in the floor-corrected bump runs, experiment \code{bump\_floor}). The online budget is varied over $10\,000$, $30\,000$, $80\,000$ and $100\,000$ particles, with every fit on the whole online cloud: the $100\,000$ column is experiment \code{suite\_lagged}, the other three are \code{suite\_budget\_tuned} (\code{scripts/budget\_sweep.py}). The same four heads on the short-trained explicit body and on the implicit body are also in the store (experiment \code{suite\_lagged}) and are summarised in the study log (\code{paper/studies\_log.md}).

\begin{table}[h]\centering\footnotesize
\caption{Frozen tuned or searched body ($500\,000$ particles) with an online head, against NW re-solved from scratch, by online particle budget: pooled MAE in vol bp, $20$ SSVI scenarios, two online seeds, both lags; liquid MAE at $100\,000$; online seconds at $10\,000$ and $100\,000$ particles (median). Bold: the best method row in each column; the PDE floor is the reference and is not a contender, and a zero online time (no online work) does not count.}\label{tab:suite-heads}
\resizebox{\textwidth}{!}{\begin{tabular}{l rrrrr rrrrr rr}
\toprule
method & \multicolumn{5}{c}{surface lag} & \multicolumn{5}{c}{surface plus $2\%$ spot lag} & \multicolumn{2}{c}{online s}\\
\cmidrule(lr){2-6}\cmidrule(lr){7-11}\cmidrule(lr){12-13}
& 10k & 30k & 80k & 100k & liquid & 10k & 30k & 80k & 100k & liquid & 10k & 100k\\
\midrule
NW re-solve on $S_1$ & -- & -- & -- & 44 & -- & -- & -- & -- & 44 & -- & -- & 8.5\\
Tuned body, stale $f$, fresh Dupire & -- & -- & -- & -- & -- & -- & -- & -- & -- & -- & -- & --\\
Tuned body + spline head & -- & -- & -- & -- & -- & -- & -- & -- & -- & -- & -- & --\\
Tuned body + RKHS head & -- & -- & -- & -- & -- & -- & -- & -- & -- & -- & -- & --\\
Tuned body + ridge head & -- & -- & -- & -- & -- & -- & -- & -- & -- & -- & -- & --\\
\bottomrule
\end{tabular}
}
\end{table}

\emph{Results} (\cref{tab:suite-heads}). A frozen tuned body with a guarded online head is at least as good as a fresh NW calibration at every online budget, and better where the budget is small. At $10\,000$ online particles NW re-solved from scratch is at $50\bp$ under both lags; the spline head is at $46$, the RKHS and ridge heads at $46$, and the advantage shrinks as the online cloud grows: at $80\,000$ particles, the practical budget, the spline head is $43$ against NW's $44$--$45$ under both lags, at the same online time ($7$ s), and the ridge head is the best row after the spot move ($43$) at three times the time; at $100\,000$ every head is within $1$--$2\bp$ of NW and $1\bp$ of each other. On liquid quotes the heads are at $15$--$16\bp$ against NW's $17$--$18$ at every budget. The head is doing the work: the stale body with a fresh Dupire numerator is $50\bp$ under the surface lag and $88$ after the spot move, at every budget, because the stored $\E[V\mid X]$ was fitted for a cloud started at the old spot. The spline head is the row to use: the cheapest head, never worse than the others, and it turns the tuned body's cold-suite draw into a small but consistent win over the re-solve at every budget from $10$k to $100$k. The searched body with the spline head is the best row at every budget from $30\,000$ up: $42\bp$ under the surface lag at $80\,000$ and $100\,000$ against NW's $44$, and $43$ after the spot move against $45$ and $44$, with the liquid-quote error at $16\bp$ against NW's $17$--$18$, and the fastest online sweep ($7.1$ s at $100\,000$ against $8.5$ for NW), because the network is half the width. Its liquid-quote deficit as a cold calibrator disappears once the head runs on it. At $10\,000$ the searched and tuned bodies tie ($46\bp$). The search bought about half a basis point over the hand-tuned body at the practical budget and a faster head; the two-basis-point margin over the re-solve is the body-plus-head design, not the search.


\subsection{Reproducibility}\label{sec:suite-repro}

Each run records the git hash, the full configuration as parameters, the leverage field, the implied-volatility error grid, and the trained objects. A loader (\code{nparticle suite load <run\_id>}) rebuilds the leverage field and, where present, the network or per-slice estimators from the artifacts alone, so any row of the tables can be re-scored or reused as an offline body without recomputation.

\section{Importance sampling}\label{sec:tilt}

The particle method estimates $\E[V_t\mid X_t]$ from wherever the particles happen to be, and at short maturities almost none of them are in the wings. A volatility-preserving Girsanov tilt moves them there: the cloud is simulated as a defensive mixture of three drifts $-\theta,0,+\theta$ on log-spot with weights $\tfrac14,\tfrac12,\tfrac14$, the exact discrete-time weights are accumulated along each path and bounded by $2$, and every per-slice fit becomes a weighted regression (\code{calibrate/importance.py}). The online cost is one drift term and three running scalars per particle. This section measures what the tilt buys, where, and what stops it buying more.

The summary is this. For the classical kernel estimator the tilt removes a fifth to a third of the wing variance at every particle budget and lowers the leverage error against the PDE reference by $12$ to $14$ percent, so that $30\,000$ tilted particles match $80\,000$ untilted ones on both. On the quoted surface that is worth $3\bp$ at $80$k. On far-wing quotes, where the tilt changes the leverage, it is worth $4\bp$ for the kernel method and $2.4\bp$ for the frozen body with an online head, at no online cost. The network estimator does not benefit until its loss is reweighted by the sampling density, after which it does.

A design is a schedule for $\theta$ and a cap $c$ on $\sum_k\theta_k^2\,\Delta t/(1-\rho^2)$, the log-weight variance of a pure tilt. A constant schedule has $\theta=\sqrt{c(1-\rho^2)/T}$; an inverse-square-root schedule, $\theta_k=a/\sqrt{\max(t_k,\Delta t)}$ with $a$ fixed by the same cost, concentrates the tilt at short times; a front-loaded schedule spends the whole cost on $[0,0.25)$ and then stops. The sweep is $c\in\{1,3,9\}$ for the first two schedules and $c\in\{3,8\}$ for the third, plus untilted.

\subsection{Slice layer}\label{sec:tilt-slices}
The tilt acts on estimation variance, so it is measured first on its own, without the feedback of the calibration loop. Particles are simulated under the frozen PDE-reference leverage of each of the six study scenarios (s01, s02, s05, s09, s11, s16), and at the four quoted maturities NW and the searched network are fitted on the weighted cloud and evaluated at the $13$ quoted strikes against the PDE conditional expectation, at $10$k, $30$k and $80$k particles with five seeds. The error at each strike splits into its standard deviation over seeds, which the tilt can remove, and its mean, the shared Euler-versus-Fokker--Planck bias, which it cannot. NW's bandwidth is Silverman's rule on the weighted spread of the cloud, so both arms smooth on the same scale. A design wins on the smallest NW wing standard deviation at $30$k, and is disqualified if the bias it adds exceeds one percentage point plus half the standard deviation it removes.

\begin{table}[h]\centering\footnotesize
\caption{Slice layer: wing standard deviation over seeds (at $T=0.25$ and pooled over maturities) and wing bias, in percent of $\E[V\mid X]$, by estimator and budget, untilted and under the constant designs of cost $3$ and $9$ and the front-loaded design of cost $8$; the inverse-square-root designs and the other costs are in experiment \code{tilt\_slices}. Bold: smallest in each standard-deviation column; the bias columns are signed and are not bolded.}\label{tab:tilt-slices}
\resizebox{\textwidth}{!}{\begin{tabular}{l rrr rrr rrr rrr}
\toprule
estimator, $N$ & \multicolumn{3}{c}{none} & \multicolumn{3}{c}{constant-3} & \multicolumn{3}{c}{constant-9} & \multicolumn{3}{c}{front-8}\\
& std $T{=}0.25$ & std & bias & std $T{=}0.25$ & std & bias & std $T{=}0.25$ & std & bias & std $T{=}0.25$ & std & bias\\
\midrule
NW, 10k & 25.5 & 13.6 & -8.2 & 23.6 & 11.7 & -7.9 & 18.7 & 10.7 & -7.4 & 14.8 & 11.2 & -7.0\\
NW, 30k & 19.3 & 10.4 & -6.2 & 19.6 & 9.3 & -6.5 & 14.4 & 8.0 & -6.3 & 11.2 & 8.6 & -6.1\\
NW, 80k & 20.0 & 8.8 & -4.6 & 11.7 & 6.5 & -7.1 & 11.6 & 6.2 & -6.1 & 7.7 & 6.3 & -5.8\\
Net, 10k & 11.6 & 8.1 & -9.8 & 7.5 & 7.4 & -9.7 & 7.5 & 7.5 & -10.5 & 7.9 & 8.1 & -10.9\\
Net, 30k & \textbf{5.4} & \textbf{5.1} & -9.5 & \textbf{5.3} & 5.7 & -9.9 & \textbf{4.1} & 5.5 & -10.1 & 3.9 & 5.8 & -12.2\\
Net, 80k & 8.2 & 6.3 & -8.9 & 6.0 & \textbf{5.2} & -10.2 & 6.1 & \textbf{5.3} & -11.0 & \textbf{3.3} & \textbf{4.6} & -11.8\\
\bottomrule
\end{tabular}
% winner: constant-9
}
\end{table}

\begin{figure}[h]\centering
\includegraphics[width=\textwidth]{../figures/out/tilt_cloud_s02.png}
\includegraphics[width=\textwidth]{../figures/out/tilt_cloud_s11.png}
\caption{Particle clouds at $30$k under the frozen PDE leverage, untilted (left) and under the winning design (right), at $t=0.25$ (top) and $t=1$ (bottom), with the PDE conditional expectation (solid), the NW (dashed) and network (dotted) estimates, and the quoted wing strikes; marker area is the importance weight. s02 (Feller satisfied) and s11 (Feller violated).}\label{fig:tilt-cloud}
\end{figure}

\begin{figure}[h]\centering
\includegraphics[width=\textwidth]{../figures/out/tilt_leverage_s02.png}
\includegraphics[width=\textwidth]{../figures/out/tilt_leverage_s11.png}
\caption{The same slices in leverage: PDE (solid) against the NW and network estimates, with the weighted marginal density of the cloud behind.}\label{fig:tilt-leverage}
\end{figure}

The winner is the strongest constant tilt, $c=9$ (\cref{tab:tilt-slices}, \cref{fig:tilt-cloud,fig:tilt-leverage}). For NW it takes the pooled wing standard deviation from $13.6$, $10.4$ and $8.8$ percent at $10$k, $30$k and $80$k to $10.7$, $8.0$ and $6.2$, and the three-month figure from $25.5$, $19.3$ and $20.0$ to $18.7$, $14.4$ and $11.6$. Tilted $30$k ($8.0$) sits at untilted $80$k ($8.8$), which is the budget equivalence the tilt was built for. The bias is untouched, $-6$ to $-8$ percent at every design. The front-loaded schedule reaches the three-month wings better still (NW's three-month standard deviation falls to $14.8$, $11.2$ and $7.7$, and the far right strike goes from $13$ effective particles to $102$) at the same pooled figure, but its frozen sample fraction of $0.6$ costs every later slice end to end, so the constant design is the one carried forward.

\subsection{End to end}\label{sec:tilt-e2e}
Cold: NW on the $20$ SSVI scenarios and the searched network on the six study scenarios, at $10$k, $30$k and $80$k particles, seeds $0$ and $1$, untilted and under the winning design, scored as \cref{sec:suite-cold} (experiment \code{tilt\_cold}). Online: the searched $500\,000$-particle body with the spline head on a tilted online cloud at $10$k and $80$k, both lags (experiment \code{tilt\_online}), against the untilted cells of \cref{tab:suite-heads}; the P-spline is fitted by weighted least squares. Repricing is untilted throughout.

\begin{table}[h]\centering\footnotesize
\caption{Cold calibration with and without the tilt by budget: pooled, wings, three-month and liquid MAE in vol bp, leverage RMSE against the PDE reference, the minimum per-slice effective sample fraction, and median seconds. NW rows are over the $20$ SSVI scenarios; the network rows are over the six study scenarios, as their labels say, and the floor row is over the $20$. Bold: smallest in each column over all rows but the floor, so a bold tilted row at a small budget reads as budget equivalence.}\label{tab:tilt-cold}
\resizebox{\textwidth}{!}{\begin{tabular}{l rrrr r r r}
\toprule
method, $N$, arm & pooled & wings & $T{=}0.25$ & liquid & lev.\ RMSE & min ESS & s\\
\midrule
NW (20 scenarios), 10k, untilted & 51 & 74 & 107 & 21 & 0.124 & -- & \textbf{0.7}\\
NW (20 scenarios), 10k, tilted & 50 & 73 & 107 & 20 & 0.108 & 0.58 & \textbf{0.7}\\
NW (20 scenarios), 30k, untilted & 45 & 66 & 98 & 18 & 0.102 & -- & 2.4\\
NW (20 scenarios), 30k, tilted & 47 & 70 & 107 & 17 & 0.089 & 0.58 & 2.7\\
NW (20 scenarios), 80k, untilted & 46 & 69 & 105 & \textbf{16} & 0.092 & -- & 6.6\\
NW (20 scenarios), 80k, tilted & \textbf{43} & 65 & 101 & \textbf{16} & \textbf{0.082} & 0.58 & 7.5\\
Explicit NN, searched (six study scenarios), 10k, untilted & 58 & 83 & 112 & 22 & 0.252 & -- & 937.3\\
Explicit NN, searched (six study scenarios), 10k, tilted & 60 & 85 & 115 & 23 & 0.216 & 0.58 & 947.0\\
Explicit NN, searched (six study scenarios), 30k, untilted & 53 & 76 & 106 & 21 & 0.188 & -- & 949.7\\
Explicit NN, searched (six study scenarios), 30k, tilted & 59 & 85 & 114 & 20 & 0.204 & 0.58 & 955.2\\
Explicit NN, searched (six study scenarios), 80k, untilted & 54 & 76 & 104 & 20 & 0.167 & -- & 975.1\\
Explicit NN, searched (six study scenarios), 80k, tilted & 54 & 78 & 109 & 19 & 0.180 & 0.58 & 959.2\\
Explicit NN, density-weighted (six study scenarios), 10k, untilted & 57 & 79 & 112 & 25 & 0.267 & -- & 894.5\\
Explicit NN, density-weighted (six study scenarios), 10k, tilted & 51 & 70 & 106 & 22 & 0.219 & 0.58 & 896.2\\
Explicit NN, density-weighted (six study scenarios), 30k, untilted & 48 & \textbf{64} & \textbf{95} & 24 & 0.207 & -- & 901.8\\
Explicit NN, density-weighted (six study scenarios), 30k, tilted & 54 & 73 & 108 & 24 & 0.184 & 0.58 & 903.8\\
Explicit NN, density-weighted (six study scenarios), 80k, untilted & 50 & 68 & 103 & 21 & 0.157 & -- & 932.0\\
Explicit NN, density-weighted (six study scenarios), 80k, tilted & 49 & 66 & 98 & 23 & 0.168 & 0.58 & 932.5\\
PDE (attainable floor) & 42 & 62 & 110 & 13 & 0.000 & -- & 0.0\\
\bottomrule
\end{tabular}
}
\end{table}

\begin{table}[h]\centering\footnotesize
\caption{Searched body with the spline head on an untilted and a tilted online cloud, by online budget and lag: pooled, wings and liquid MAE in vol bp and online seconds.}\label{tab:tilt-online}
\begin{tabular}{l rrr r}
\toprule
budget, lag, arm & pooled & wings & liquid & online s\\
\midrule
10k, surface, untilted & 46 & 67 & 19 & \textbf{1.0}\\
10k, surface, tilted & 45 & 65 & 19 & 1.1\\
10k, surface\_spot, untilted & 46 & 67 & 19 & \textbf{1.0}\\
10k, surface\_spot, tilted & 47 & 69 & 20 & 1.1\\
80k, surface, untilted & \textbf{42} & 63 & \textbf{16} & 6.1\\
80k, surface, tilted & \textbf{42} & \textbf{62} & \textbf{16} & 6.7\\
80k, surface\_spot, untilted & 43 & 64 & \textbf{16} & 6.1\\
80k, surface\_spot, tilted & 43 & 64 & \textbf{16} & 6.6\\
\bottomrule
\end{tabular}

\end{table}

For NW the tilt lowers the leverage RMSE by $12$ to $14$ percent at every budget ($0.124$, $0.102$, $0.092\to0.108$, $0.089$, $0.082$ at $10$k, $30$k, $80$k), better in $39$ of $40$ scenario-seed pairs at $80$k, and tilted $30$k again matches untilted $80$k. On the quoted surface it is worth $3\bp$ at $80$k ($46\to43\bp$ pooled, $69\to65$ in the wings, better in $32$ of $40$ pairs), is within noise at $10$k, and is $2\bp$ worse at $30$k, where the loss sits at three months ($98\to107\bp$, $t=3.6$) while the leverage error still improves; the mechanism behind that last point is not settled. The frozen body with the spline head is unchanged at every budget and lag ($42$--$43\bp$ at $80$k either way): that row already sits on the floor of the quoted surface. The plain searched network gains nothing; \cref{sec:tilt-dw} says why and fixes it.

Two things cap what the tilt can do on the quoted surface, and both are measurable from the stored runs. The three-month right side is bias, not variance: at $k=0.14$ to $0.31$ the slice-layer error in $\E[V\mid X]$ is $-16$ to $-19$ percent for every design, with a standard deviation of $1$ to $6$, the Euler cloud disagreeing with the reference where $V$ sits at its zero atom. And the scoring floor at $200$ steps is mostly time-discretisation bias. Over the $20$ scenarios the PDE-exact leverage repriced at $400$ steps scores $28.4\bp$ pooled, $43.2$ in the wings, $63.8$ at three months and $8.5$ on liquid quotes, against $42.0$, $61.9$, $110.3$ and $13.4$ at $200$ steps; NW itself, calibrated and repriced at $400$ steps, goes from $45.6$ to $34.3\bp$ pooled at $80$k (experiments \code{suite\_pde\_floor\_s400}, \code{tilt\_cold\_s400}). On the six hardest scenarios NW at $200$ steps even sits below the ``floor'' ($52$ against $55$), because a method calibrated on the discrete scheme compensates part of the scheme's bias while the PDE-exact leverage cannot; at $400$ steps the floor is a floor everywhere. The tilt survives the finer step: at $400$ steps it still lowers NW's leverage error in $39$ of $40$ pairs and the pooled score by $1.4\bp$ ($t=-3.3$). The time step, not the estimator, is the lever on the absolute numbers of \cref{sec:suite}, and every stored body interpolates in time, so section 4 can be re-scored at $400$ steps without retraining.

\subsection{Far-wing quotes}\label{sec:tilt-far}
The quoted grid stops at $|k|\approx0.5$ and its three-month wings fail inversion, so the score of \cref{sec:tilt-e2e} never sees the region where the tilt changes the leverage. Here every stored field is re-scored, without recalibration, on $21$ log-strikes from $0.45$ to $2.2$ spot ($|k|$ up to $0.8$) with an importance-sampled out-of-the-money repricer: puts below the money, calls above, on a cloud tilted by the constant design of cost $3$ with exact discrete-time weights, inverted through put-call parity. A quote whose target price is below $10^{-5}$ of spot is not scored, which removes the three-month far wings. The far-wing error is the mean absolute implied-vol error over the scored quotes with $|k|>0.25$.

\begin{table}[h]\centering\footnotesize
\caption{Far-wing implied-vol MAE in bp at the $80$k budget: pooled over maturities and at $0.5$, $1$ and $2$ years, the mean number of scored far-wing quotes, and the $13$-strike pooled MAE for reference. Every row is over the same $20$ SSVI scenarios and two seeds. Bold: smallest in each column over the non-floor rows.}\label{tab:tilt-far}
\begin{tabular}{l rrrrrr}
\toprule
row & far wings & $T{=}0.5$ & $T{=}1$ & $T{=}2$ & quotes & pooled (13-strike)\\
\midrule
NW, 200 steps, untilted & 37 & 46 & 35 & 19 & 40 & 46\\
NW, 200 steps, tilted & \textbf{32} & \textbf{40} & 31 & 18 & 40 & 43\\
Explicit NN, searched, untilted & 38 & 48 & 35 & 20 & 40 & 48\\
Explicit NN, searched, tilted & 38 & 47 & 36 & 20 & 40 & 49\\
Explicit NN, density-weighted, untilted & 34 & 47 & 34 & 18 & 40 & 45\\
Explicit NN, density-weighted, tilted & 33 & 44 & 30 & 18 & 40 & 44\\
NW re-solve (online) & 37 & 47 & 35 & 20 & 42 & 44\\
Searched body + spline head, untilted & 34 & 44 & 29 & 18 & 42 & \textbf{42}\\
Searched body + spline head, tilted & \textbf{32} & 43 & \textbf{26} & \textbf{15} & 42 & \textbf{42}\\
PDE floor, 200 steps & 21 & 29 & 15 & 8 & 40 & 42\\
\bottomrule
\end{tabular}

\end{table}

This is where the tilt pays (\cref{tab:tilt-far}). For NW at $80$k the far-wing error falls from $36.6$ to $32.5\bp$, $-4.1\bp$ paired over $40$ scenario-seed pairs ($t=-5.9$, better in $34$), at every maturity ($45.7\to39.6$ at six months, $35.4\to31.0$ at one year, $18.9\to17.6$ at two), a quarter of its gap to the $200$-step floor of $20.7$. The frozen searched body with the spline head on a tilted online cloud improves by $2.4\bp$ under both lags ($t=-3.9$ and $-3.6$, better in $30$ of $40$), and that row beats the NW re-solve at the same budget by $5.5$ to $5.9\bp$ ($t=-6.9$ and $-8.1$): the offline body carries the wings and the tilted online cloud corrects them. At $400$ steps NW's far-wing error is $25.2\bp$ untilted, $11\bp$ from the time step alone, and the tilt is still worth $1.5\bp$ on top ($t=-2.2$, better in $28$ of $40$). The plain searched network sits at NW's untilted level ($37.7$) and does not move.

\subsection{The network and the density-weighted fit}\label{sec:tilt-dw}
On s11 at three months the searched network's $\E[V\mid X]$ is exactly zero from $k=0.14$ outward, so its leverage sits on the cap $\Lmax=4$ across the right wing (\cref{fig:tilt-leverage}), while NW on the same cloud gives $0.007$ to $0.056$ against the PDE's $0.012$ to $0.064$. No prior is responsible; switching the monotone penalty, the mean matching and the tail rule off changes nothing. The cause is the loss. About $900$ of $80\,000$ particles sit at $k>0.2$, and least squares gives that wing an influence equal to its probability mass, one percent, while the bulk, where $f$ falls steeply to zero near $k=0.06$, sets the shape; the network extrapolates the descent and its softplus output clamps at zero. NW is local and never sees the bulk from the wing. The same fact is why the tilt cannot help the network as built: it puts $1475$ particles in that region instead of $906$, each with an importance weight below one, and the weighted loss leaves the wing's total weight at its physical mass, whereas NW's local ratio gains from the particle count directly.

The fix is to weight the loss by $w(x)/\hat p(x)$, the importance weight divided by a kernel estimate of the physical density of $X$ at the particle, floored at $1/50$ of the modal density and normalised to mean one (\code{NNRegressor(design\_weight=True)}, cold row \code{explicit\_nn\_opt\_dw}, the promoted recipe otherwise unchanged). Any positive weight that depends on $x$ alone leaves the minimiser at $\E[V\mid X]$, so the fit stays unbiased and the wings count by their particles. On the s11 slice the network then recovers the whole right wing ($0.008$, $0.020$, $0.032$, $0.040$, $0.043$ at the five outer strikes), smoother than NW. In the slice layer at $80$k its three-month wing bias halves ($-14.6$ to $-7.0$ percent) with a third of NW's three-month variance ($6.3$ against $20.0$), at the cost of pooled wing variance ($6.3$ to $7.2$), and the tilt now reaches it (three-month standard deviation $6.3\to4.7$).

End to end at $80$k on the $20$ scenarios (\cref{tab:tilt-cold,tab:tilt-far}, $40$ pairs), the density-weighted network is $3.6\bp$ better than the plain one pooled ($t=-2.9$), $7.9$ in the wings ($t=-3.1$) and $3.3$ in the far wings ($t=-2.2$). Against NW it is a draw pooled ($44.6$ against $45.6$), $5.6\bp$ better in the wings ($63.4$ against $69.0$, $t=-1.9$), $2.2$ better in the far wings ($t=-2.0$), and $3.8\bp$ worse on liquid quotes ($19.7$ against $15.9$, $t=3.4$). Equalising the loss across log-spot buys the wings with near-the-money accuracy, and the cap on the density weight, set at one value here, is the dial between the two. Under the tilt the density-weighted network is the best cold row in the wings ($60.4\bp$ against NW's $65.0$), ties NW pooled ($43.6$ against $43.4$) and in the far wings ($32.5$ each), and the tilt now moves it in the far wings by $-1.9\bp$ ($t=-2.7$). At $30$k on the six study scenarios it is again the best untilted cold row ($48.1\bp$ against NW's $52.0$ and the plain network's $53.3$). What remains of its deficit to NW is in the bulk (leverage RMSE $0.187$ against $0.149$), and the offline body of \cref{sec:suite-heads} has not yet been retrained with this loss.

\end{document}
